\chapter{Stability--Optimality Duality}
\label{ch:duality}

A useful swing should achieve its impact objective with tolerable effort and sensitivity to error. These requirements are related, but none implies the others without a model and a stated criterion. Linear quadratic control gives a particularly clear connection: the same Riccati solution describes the minimum future quadratic cost and supplies a certificate of closed-loop stability. This is a mathematical identity under declared assumptions, not evidence that a golfer solves a Riccati equation or minimizes mechanical energy.

This chapter develops that identity in continuous and discrete time, checks its rates and numerical predictions, and separates it from several tempting overextensions. An optimal-control cost need not be a physical energy; an optimal state path need not be a geodesic; a well-conditioned Riccati matrix need not imply controllability; and an observer-based implementation need not retain full-state feedback robustness. These distinctions sharpen the golf question: which perturbations can the available inputs correct, which can the measurements reveal, and at what cost within the remaining time to impact?

\section{What the Quadratic Problem Assumes}

\subsection{State, Port and Objective}

Consider the continuous-time system
\begin{equation}
\dot x=Ax+Bu,\qquad
J(x_0,u)=\int_0^\infty (x^TQx+u^TRu)\,dt,
\end{equation}
where $Q=Q^T\succeq0$ and $R=R^T\succ0$. Here $x$ is a deviation from an equilibrium with the equilibrium input removed. For a swing trajectory the analogous variables are deviations from a feasible time-dependent reference; that problem generally has time-dependent matrices and a finite horizon. The actuator port matters: torque, muscle excitation and motor voltage lead to different $B$, retained actuator states and admissible inputs.

The components of $x$ and $u$ require explicit units or reference scales. If $x=Tz$ and $u=Uv$ for invertible scale matrices, the transformed data are
\begin{equation}
A_z=T^{-1}AT,\quad B_z=T^{-1}BU,\quad
Q_z=T^TQT,\quad R_z=U^TRU.
\end{equation}
Consequently $S_z=T^TST$ and $K_z=U^{-1}KT$. The physical policy is unchanged by this consistent conversion. Raw eigenvalues, Euclidean norms and condition numbers of $S$ generally change. Comparing an angle gain with an angular-velocity gain by their numerical magnitudes alone compares quantities with different units.

The quadratic criterion is a choice about performance. Penalizing activation squared can be an effort surrogate; penalizing torque squared is a different surrogate. Neither is automatically metabolic expenditure, joint work, tissue loading or shaft energy. For example, joint mechanical power is torque times angular velocity, with a sign and an energy-transfer interpretation absent from a positive quadratic torque penalty. State exactly which criterion an argument uses.

\subsection{Existence and Positive Definiteness}

For the infinite-horizon continuous-time LQR problem, standard sufficient hypotheses are stabilizability of $(A,B)$ and detectability of $(Q^{1/2},A)$. Stabilizability means that modes which cannot be controlled are already asymptotically stable. Detectability means that modes invisible to the state penalty are already asymptotically stable. Under these hypotheses the stabilizing algebraic Riccati solution is symmetric positive semidefinite, the feedback stabilizes the system, and the optimal value is $x_0^TSx_0$.

Positive semidefinite is not positive definite. A stable unpenalized state can have zero value, so $S$ can be singular even when the feedback is stabilizing. Assuming $Q\succ0$ is a convenient sufficient condition for $S\succ0$ in the setting here. More generally an appropriate observability condition on the penalized state excludes zero-cost directions. We will state $S\succ0$ explicitly whenever it defines a norm or metric. Detectability alone does not justify an inverse $S^{-1}$.

The corresponding discrete-time hypotheses replace asymptotic stability by the condition that eigenvalues lie strictly inside the unit circle. These assumptions concern the declared mathematical model. They do not verify that a linearization remains accurate through a large swing, actuator saturation or a contact transition.

\section{The Continuous-Time Identity}

\subsection{Completing the Square}

Let $S=S^T$ solve
\begin{equation}
A^TS+SA-SBR^{-1}B^TS+Q=0,
\qquad K=R^{-1}B^TS.
\end{equation}
For $V(x)=x^TSx$, direct differentiation gives
\begin{equation}
\begin{aligned}
\dot V&=x^T(A^TS+SA)x+2x^TSBu,\\
x^TQx+u^TRu+\dot V&=(u+Kx)^TR(u+Kx).
\end{aligned}
\end{equation}
Every sign is consequential. Integrating the second identity over $[0,T]$ gives
\begin{equation}
J_T=V(x_0)-V(x(T))+
\int_0^T(u+Kx)^TR(u+Kx)\,dt.
\end{equation}
For the stabilizing solution and admissible infinite-horizon comparison controls with the required vanishing terminal value, this proves optimality of $u=-Kx$ and $J^*=V(x_0)$. The terminal term cannot simply be discarded for an arbitrary nondecaying trajectory. The usual LQR existence theorem supplies the admissibility and limiting argument under the preceding hypotheses.

On the optimal closed loop $A_c=A-BK$, define $D=Q+K^TRK$. Then
\begin{equation}
A_c^TS+SA_c=-D,\qquad \dot V=-x^TDx.
\end{equation}
The feedback effort term is part of the dissipation identity. Dropping it changes the proof and can change a claimed rate. If $S\succ0$ and $D\succ0$, this is a strict quadratic Lyapunov certificate. With $D\succeq0$, a nonpositive derivative alone is insufficient to prove asymptotic convergence; an invariance or detectability argument is needed to exclude nondecaying zero-dissipation motions.

\subsection{A Value Rate Is Twice a Norm Rate}

Suppose $S\succ0$ and $D\succ0$. Define
\begin{equation}
\beta=\lambda_{\min}(S^{-1/2}DS^{-1/2})>0.
\end{equation}
Since $x^TDx\ge\beta x^TSx$, integration yields
\begin{equation}
V(t)\le e^{-\beta t}V(0),\qquad
\|x(t)\|_S\le e^{-\beta t/2}\|x(0)\|_S,
\end{equation}
where $\|x\|_S=\sqrt{x^TSx}$. In a fixed scaled Euclidean norm,
\begin{equation}
\|x(t)\|_2\le\sqrt{\kappa_2(S)}\,
e^{-\beta t/2}\|x(0)\|_2.
\end{equation}
A simpler, potentially more conservative, value rate is $\lambda_{\min}(D)/\lambda_{\max}(S)$. The maximum eigenvalue of $S$, rather than its minimum, belongs in this denominator. The factor of one half in the norm exponent follows from taking a square root. These are certified bounds, not necessarily the slowest pole or the observed decay along a particular initial direction.

For the scalar system $\dot x=x+u$ with $Q=R=1$, the stabilizing solution is $S=1+\sqrt2$ and $A_c=-\sqrt2$. Hence $x(t)=e^{-\sqrt2t}x_0$, while $V(t)=e^{-2\sqrt2t}V_0$. Calling $2\sqrt2$ the state-norm decay rate would be wrong even in this one-dimensional example. At $t=0.7$, the state multiplier is approximately $0.371595$.

\section{The Discrete-Time Identity}

\subsection{Bellman Completion and the Riccati Equation}

For $x_{k+1}=Ax_k+Bu_k$ with cost $\sum_{k=0}^\infty(x_k^TQx_k+u_k^TRu_k)$, set
\begin{equation}
H=R+B^TSB,\qquad K=H^{-1}B^TSA.
\end{equation}
The discrete algebraic Riccati equation is
\begin{equation}
S=Q+A^TSA-A^TSB(R+B^TSB)^{-1}B^TSA.
\end{equation}
It differs from the continuous-time algebraic equation. A fixed point of the discrete backward Riccati recursion solves this discrete equation, not its continuous counterpart. Expanding the quadratic terms gives
\begin{equation}
\begin{aligned}
&x^TQx+u^TRu+(Ax+Bu)^TS(Ax+Bu)-x^TSx\\
&\hspace{2em}=(u+Kx)^TH(u+Kx).
\end{aligned}
\end{equation}
On $u=-Kx$, therefore,
\begin{equation}
A_c^TSA_c-S=-Q-K^TRK=-D,
\qquad V_{k+1}-V_k=-x_k^TDx_k.
\end{equation}
The $K^TRK$ term follows from expanding both next-state value and input cost; it cannot be added after using an inconsistent expression for $A_c^TSA_c$.

\subsection{Discrete Norm and Value Factors}

With $S\succ0$ and $D\succ0$, define $\delta=\lambda_{\min}(S^{-1/2}DS^{-1/2})$ and $\rho=1-\delta$. The identity $S-D=A_c^TSA_c\succeq0$ ensures $0<\delta\le1$. It follows that
\begin{equation}
V_k\le\rho^k V_0,\qquad
\|x_k\|_2\le\sqrt{\kappa_2(S)}\,\rho^{k/2}\|x_0\|_2.
\end{equation}
If $\rho=0$, the identity implies $A_c=0$ and the state reaches zero after one step. Otherwise the norm multiplier per step is $\sqrt\rho$, not $\rho$. Replacing $\delta$ by $\lambda_{\min}(D)/\lambda_{\max}(S)$ gives another valid, generally looser, bound. A physical decay rate per second also requires the sample period; a discrete multiplier alone has no time unit.

For $A=1.2$ and $B=Q=R=1$, the scalar closed-loop multiplier is approximately $0.406472$. Its squared value, approximately $0.165219$, is the value multiplier. This elementary check catches the same norm/value confusion as the continuous-time example.

\section{Finite Horizons and Moving References}

\subsection{The Terminal Condition Matters}

For a feasible nominal trajectory, the first-order deviation dynamics are $\dot\xi=A(t)\xi+B(t)v$. With terminal penalty $\xi(T)^TS_f\xi(T)$, the local quadratic value is $V(t,\xi)=\xi^TS(t)\xi$ and
\begin{equation}
-\dot S=A^TS+SA-SBR^{-1}B^TS+Q,\qquad S(T)=S_f.
\end{equation}
Here $K(t)=R(t)^{-1}B(t)^TS(t)$. Including $\partial_t V$ restores the same completion identity. Along the local optimal feedback,
\begin{equation}
\frac{d}{dt}(\xi^TS\xi)=-\xi^T(Q+K^TRK)\xi.
\end{equation}
This is the chosen homogeneous quadratic tracking problem in the deviation variables. Expanding a general original cost about a feasible but nonoptimal trajectory can also produce linear terms in the local value and a feedforward correction. Feasibility alone does not make those terms vanish or make the reference locally optimal.
In discrete time use $S_{k+1}$ in $H_k=R_k+B_k^TS_{k+1}B_k$ and $K_k=H_k^{-1}B_k^TS_{k+1}A_k$, then propagate
\begin{equation}
S_k=Q_k+A_k^TS_{k+1}A_k-
A_k^TS_{k+1}B_kH_k^{-1}B_k^TS_{k+1}A_k.
\end{equation}
Both formulas run backward from the terminal condition. A finite-horizon optimum does not by itself establish infinite-time stability, recursive feasibility or stability under repeated replanning. With $S_f=0$, the value vanishes at $T$ for every terminal state. It is then impossible to regard $S(T)$ as a positive-definite metric or infer a nonzero terminal error bound from terminal value alone.

Uniform bounds $mI\preceq S(t)\preceq MI$ and a uniform strict differential inequality are needed for a uniform exponential norm estimate. For a nonlinear swing, the local feedback and quadratic value are approximations near the reference. Higher-order remainders, input limits, estimation error and mode transitions must be controlled on a stated neighborhood before extending the certificate to the nonlinear motion.

\subsection{Impact Objectives and Remaining Control Authority}

An impact objective may penalize predicted clubhead velocity, face orientation and strike location rather than all joint coordinates equally. If $y=H(x)$ is the impact output and $C_T=DH(\bar x(T))$, a local terminal penalty takes the form $S_f=C_T^TWC_T$. This matrix is often semidefinite because multiple state changes produce the same first-order impact output. Those null directions are not automatically harmless: they may affect constraints, higher-order outcomes or subsequent dynamics.

This expression is exact for the squared linearized output error. For the nonlinear penalty $(H(x)-y_*)^TW(H(x)-y_*)$, its half-Hessian also contains weighted second derivatives of $H$ multiplied by the nominal output residual. Those additional terms vanish when $H(\bar x(T))=y_*$. Using only $C_T^TWC_T$ at a nonzero residual is a Gauss--Newton approximation, not the full terminal curvature.

The remaining horizon and the state transition matrix determine how a torque change can affect that output. The Riccati solution combines that dynamics with chosen costs; it does not substitute for a controllability calculation. A large gain near impact may demand unavailable torque or bandwidth. For a ball-contact event whose time changes with the trajectory, output sensitivity must include the event-time correction and any reset derivative. A smooth fixed-time terminal model cannot silently stand in for this hybrid calculation.

\section{A Reproducible Discrete Example}

\subsection{Model and Numerical Authority}

Use the declared illustrative matrices
\begin{equation}
A=\begin{bmatrix}1&0.1\\0.981&1\end{bmatrix},\quad
B=\begin{bmatrix}0.005\\0.1\end{bmatrix},\quad
Q=\begin{bmatrix}10&0\\0&1\end{bmatrix},\quad R=0.1.
\end{equation}
These historical pendulum-like coefficients mix approximation orders: the state matrix resembles a forward-Euler upright-pendulum step, while the position input coefficient retains a second-order term. We treat them as the definition of a discrete example, not an exact zero-order-hold discretization of a continuous pendulum. For a physical application, derive $A_d$ and $B_d$ together from the same declared continuous model and hold assumption, then verify the approximation at the chosen sample period.

The coordinates below are scaled angle and angular-rate deviations and the input is a scaled command. A physical interpretation requires specifying those reference scales and the torque or acceleration port. The normalized initial direction $(1,0)^T$ is used to display linear homogeneity; it does not validate a one-radian upright-pendulum perturbation or a finite golf-swing intervention. Multiplying that initial direction by $\varepsilon$ multiplies the state and input by $\varepsilon$ and the value by $\varepsilon^2$.

Both editions obtain their matrices and table from the same generator, with a separate SciPy Riccati solution and trajectory propagation used in the regression checks. Displayed values are rounded only after calculation; verifying a small residual from the rounded matrices alone is inappropriate.

% DO NOT EDIT. Generated by scripts/generate_worked_examples.py.
% Every number below is solved, not transcribed; CI fails if this file
% drifts from what the generator produces (issue #3518).


\newcommand{\chsixRiccati}{94.9386 & 20.7112 \\ 20.7112 & 7.3663}
\newcommand{\chsixRiccatiEigMax}{99.5898}
\newcommand{\chsixRiccatiEigMin}{2.7150}
\newcommand{\chsixCondition}{36.6808}
\newcommand{\chsixGain}{17.1287 & 5.5643}
\newcommand{\chsixGainAngle}{17.1287}
\newcommand{\chsixGainRate}{5.5643}
\newcommand{\chsixBtS}{2.5458 & 0.8402}
\newcommand{\chsixBtSA}{3.3700 & 1.0948}
\newcommand{\chsixBtSB}{0.09675}
\newcommand{\chsixRplusBtSB}{0.19675}
\newcommand{\chsixBK}{0.0856 & 0.0278 \\ 1.7129 & 0.5564}
\newcommand{\chsixClosedLoop}{0.9144 & 0.0722 \\ -0.7319 & 0.4436}
\newcommand{\chsixEigOne}{0.6281}
\newcommand{\chsixEigTwo}{0.7298}
\newcommand{\chsixSpectralRadius}{0.7298}
\newcommand{\chsixStageCost}{39.3393 & 9.5310 \\ 9.5310 & 4.0962}
\newcommand{\chsixStageEigMax}{41.7517}
\newcommand{\chsixStageEigMin}{1.6838}
\newcommand{\chsixRho}{0.9831}
\newcommand{\chsixSqrtCondition}{6.0565}
\newcommand{\chsixRiccatiResidual}{1e-10}
\newcommand{\chsixValueRatioTen}{0.0043}
\newcommand{\chsixBoundTen}{0.8432}

\newcommand{\chsixValueTable}{%
\begin{tabular}{|c|c|c|c|c|}
\hline
Step $k$ & $V_k$ & $\norm{\dx_k}$ & $V_k/V_{k-1}$ (actual) & $\rho^k$ (bound) \\
\hline
0 & 94.9386 & 1.0000 & --- & 1.0000 \\
1 & 55.5992 & 1.1712 & 0.5856 & 0.9831 \\
2 & 33.2717 & 1.2654 & 0.5984 & 0.9665 \\
3 & 19.9314 & 1.2015 & 0.5990 & 0.9501 \\
4 & 11.8393 & 1.0561 & 0.5940 & 0.9341 \\
5 & 6.9497 & 0.8847 & 0.5870 & 0.9183 \\
6 & 4.0294 & 0.7177 & 0.5798 & 0.9027 \\
7 & 2.3094 & 0.5693 & 0.5731 & 0.8875 \\
8 & 1.3099 & 0.4441 & 0.5672 & 0.8725 \\
9 & 0.7363 & 0.3421 & 0.5621 & 0.8577 \\
10 & 0.4107 & 0.2610 & 0.5578 & 0.8432 \\
\hline
\end{tabular}%
}

\subsection{Computed Matrices and Decay Bound}
The stabilizing discrete solution and feedback are
\begin{equation}
S=\begin{bmatrix}\chsixRiccati\end{bmatrix},\qquad
K=\begin{bmatrix}\chsixGain\end{bmatrix}.
\end{equation}
The closed-loop and stage-cost matrices are
\begin{equation}
A_c=\begin{bmatrix}\chsixClosedLoop\end{bmatrix},\qquad
D=\begin{bmatrix}\chsixStageCost\end{bmatrix}.
\end{equation}
The unrounded Riccati residual is below $10^{-10}$. The closed-loop eigenvalues
are $\chsixEigOne$ and $\chsixEigTwo$, so the spectral radius is
$\chsixSpectralRadius$. The value matrix has condition number $\chsixCondition$.
Using the conservative ratio $\lambda_{\min}(D)/\lambda_{\max}(S)$ gives
\begin{equation}
\rho=\chsixRho,\qquad
\|x_k\|_2\le\chsixSqrtCondition\,\rho^{k/2}\|x_0\|_2.
\end{equation}
\subsection{Actual Trajectory Versus Its Bound}
For the normalized initial direction $x_0=(1,0)^T$, direct propagation gives:
\begin{center}
\small
\chsixValueTable
\end{center}
At step ten the actual value ratio $V_{10}/V_0$ is $\chsixValueRatioTen$,
while its conservative bound is $\chsixBoundTen$.


The Euclidean norm initially grows even though the value decreases at every step. This is compatible with a stable, nonnormal closed loop and an anisotropic value function. It refutes the claim that every norm must decrease whenever a quadratic Lyapunov function does. The bound column is $\rho^k$, a bound on $V_k/V_0$; the preceding ratio column is the actual one-step ratio $V_k/V_{k-1}$. Those quantities answer different questions and must not be compared as if both were observations of the same decay factor.

\section{Geometry Without Overclaiming}

\subsection{When a Value Defines a Metric}

For $V(x)=x^TSx$ with constant $S\succ0$, the Hessian is $2S$, not $S$. Either $S$ or $2S$ defines a positive-definite inner product after a stated constant normalization. Ellipsoids $x^TSx=c$ have axes determined by the eigenvectors of $S$, with semi-axis lengths $\sqrt{c/\lambda_i(S)}$. These are not generally eigenvectors of $A_c$. The state transition maps an initial ellipsoid to another ellipsoid, but does not generally shrink it by a common factor in every direction.

For a nonlinear value function, a Hessian need not be positive definite or even exist across changes of the optimal policy. Moreover an ordinary coordinate Hessian is not automatically a tensor. Under $x=\phi(z)$,
\begin{equation}
\frac{\partial^2 V}{\partial z_i\partial z_j}
=\frac{\partial\phi^a}{\partial z_i}
\frac{\partial\phi^b}{\partial z_j}
\frac{\partial^2 V}{\partial x^a\partial x^b}
+\frac{\partial V}{\partial x^a}
\frac{\partial^2\phi^a}{\partial z_i\partial z_j},
\end{equation}
with repeated indices summed. The second term vanishes at a stationary point; away from it, a coordinate Hessian alone does not define a coordinate-independent Riemannian metric. A covariant Hessian requires a declared connection and still needs a positivity argument. A family of local Riccati matrices along one reference similarly does not automatically define a smooth positive metric over every state and direction in a neighborhood.

The tangent bundle is the disjoint union $TM=\bigsqcup_{x\in M}T_xM$. A chart provides a local identification with a product; a global identification $TM=M\times\mathbb R^n$ requires additional structure and is not true for every manifold. When the state includes orientations or constrained configurations, this distinction prevents a convenient local coordinate representation from becoming a false global statement.

\subsection{Optimal Paths Are Not Generally Geodesics}

A Riemannian geodesic is stationary for an appropriate length or energy functional; sufficiently short regular segments locally minimize length. A geodesic need not globally minimize length or be dynamically stable. LQR minimizes a time integral of state and input penalties subject to dynamics. That functional is not generally the metric length of the state curve.

A concrete counterexample is
\begin{equation}
A=\operatorname{diag}(-1,-2),\quad B=I,\quad
Q=\operatorname{diag}(3,12),\quad R=I.
\end{equation}
The Riccati solution is $S=\operatorname{diag}(1,2)$ and $A_c=\operatorname{diag}(-2,-4)$. Starting from $(1,1)^T$, the optimal path is $(e^{-2t},e^{-4t})^T$, a parabola $x_2=x_1^2$. Geodesics of the constant metric $S$ are straight lines up to reparameterization. At $t=0$, velocity $(-2,-4)^T$ and acceleration $(4,16)^T$ are not collinear, so even a reparameterization cannot make this curve a straight line. Optimality here is exact; the geodesic interpretation fails.

The useful geometric statement is instead local and differential. For a linear closed loop, the separation $\delta x$ obeys the same dynamics as $x$, and $\delta x^TS\delta x$ satisfies the Lyapunov inequality. For a nonlinear vector field $F_c$, a contraction certificate requires a positive metric $M(x,t)$ with uniform bounds and
\begin{equation}
\partial_tM+D_xM[F_c]+(D_xF_c)^TM+M(D_xF_c)
\preceq-2\alpha M.
\end{equation}
This is a condition on the full differential dynamics throughout a stated region, together with suitable trajectory and path-domain assumptions for a finite-distance conclusion. It does not follow merely because one reference trajectory has a Riccati solution.

\subsection{Conditioning Is Not Controllability}

Take $A=-I_2$, $B=0$, $Q=I_2$ and $R=1$. There is no control authority at all, yet the stable uncontrolled dynamics have $S=I_2/2$ and $\kappa_2(S)=1$. Conversely, rescaling one state coordinate can make the condition number of a perfectly valid value matrix arbitrarily large without changing the physical closed-loop system. Thus neither a universal condition-number threshold nor a small minimum eigenvalue of $S$ is a test of controllability.

Weak control over an unstable penalized mode can instead make its value very large as input authority tends to zero. A large value can indicate costly correction under the chosen objective; it does not identify the cause without separating dynamics, input scaling and weights. Increasing $Q$ increases value in the positive-semidefinite order for the same admissible problem, but does not necessarily improve its condition number. Nor does a larger $S$ establish a larger nonlinear region of attraction: for a fixed level $c$, larger quadratic weights can make $\{x:x^TSx\le c\}$ smaller.

For a swing, inspect finite-horizon controllability or output-reachability measures with declared state and input scales, singular directions, constraints and uncertainty. Report the resulting achievable change in the impact output. The Riccati matrix then explains the chosen tradeoff among those changes rather than serving as a universal scalar measure of coordination quality.

\section{The Exact Scope of LQR Robustness}

\subsection{A Common Input-Gain Margin}

For clarity take $Q\succ0$ and the continuous-time full-state feedback above. If the actual applied feedback is $u=-\eta Kx$ for a common real scalar $\eta$, then
\begin{equation}
(A-\eta BK)^TS+S(A-\eta BK)
=-Q-(2\eta-1)K^TRK.
\end{equation}
Therefore every $\eta\ge1/2$ yields a strict Lyapunov inequality under these stronger hypotheses. The usual open interval $\eta>1/2$ is a standard guaranteed sufficient range; boundary and semidefinite-$Q$ conclusions require their own assumptions. It is not a necessary stability condition. For the scalar example $A=B=Q=R=1$, stability actually requires $\eta>1/(1+\sqrt2)\approx0.414214$. Reducing $\eta$ means weaker feedback, not more aggressive feedback.

This statement concerns a common scalar multiplication at the specified input port. It does not imply stability for every matrix gain perturbation satisfying an arbitrary norm comparison such as $\|\Delta K\|<\|K\|$. For an additive closed-loop perturbation $E$, the same certificate gives the different sufficient condition
\begin{equation}
E^TS+SE\prec D.
\end{equation}
One conservative scaled Euclidean test is $2\|S\|_2\|E\|_2<\lambda_{\min}(D)$. If the error is in the feedback gain, then $E=-B\Delta K$. Its direction and interaction with $B$ matter.

\subsection{The Return-Difference Disk Is Not a Pole Disk}

For a single input, write $B=b$, $R=r>0$ and $k=r^{-1}b^TS$. At a frequency where $j\omega I-A$ is invertible, put $z=(j\omega I-A)^{-1}b$ and $L(j\omega)=kz$. Premultiplying and postmultiplying the Riccati equation by $z^*$ and $z$, using $Az=j\omega z-b$, gives
\begin{equation}
r\bigl(|1+L(j\omega)|^2-1\bigr)=z^*Qz\ge0.
\end{equation}
Thus the Nyquist curve of the nominal return ratio avoids the open unit disk centered at $-1$. This is a frequency-response statement, not a claim that the eigenvalues of $A-BK$ lie in a disk centered at $-1$ with radius $1/2$. Continuous-time pole locations also carry inverse-time units, so a universal unscaled numerical pole disk would already need a declared time scale.

Under the usual well-posed continuous-time single-loop assumptions, the return-difference result yields the classical full-state LQR gain and phase margin guarantees. At a unity-gain crossing, $|1+e^{j\theta}|\ge1$ keeps the return ratio at least $60$ degrees from the critical negative-real direction. This familiar phase statement uses the negative-feedback return-ratio convention above. Multiple simultaneous input-channel perturbations, structured uncertainty, delays and sampled implementations require appropriately formulated robustness analysis; a SISO circle picture is not a general structured-singular-value or disk-margin certificate.

\subsection{Saturation and Estimation Change the Question}

For $\dot x=x+u$, the unconstrained LQR law is $u=-(1+\sqrt2)x$. With $|u|\le1$, at $x=2$ the saturated law gives $\dot x=1>0$, although the unconstrained law gives $\dot x=-2\sqrt2$. A local certificate can still be useful on a region where the command stays within limits, but global stability is lost in this example. Noise amplification, torque-rate limits, activation delays and contact changes likewise require an augmented model or an explicit uncertainty analysis.

When the true state is replaced by an estimate, the nominal separation principle discussed below can establish stability under its assumptions. It does not preserve all full-state LQR robustness margins. Doyle's 1978 paper \href{https://authors.library.caltech.edu/records/e39yd-evs06}{Guaranteed Margins for LQG Regulators} establishes that the general LQG class has no corresponding universal gain or phase margin guarantee. The practical conclusion is to analyze the implemented controller, estimator, sensors and actuators together.

\section{Dissipativity and Disturbance Attenuation}

\subsection{Storage Depends on the Supply Rate}

A dissipativity statement specifies a storage function $V\ge0$ and a supply rate $w$ such that $V(T)-V(0)\le\int_0^T w\,dt$. Its meaning depends on that supply rate. It is not an unconditional equivalence between dissipativity and stability. A supply that injects energy or permits state growth does not by itself establish convergence. Conversely, a Lyapunov function can supply a useful dissipativity inequality without being a physical stored energy.

The LQR completion identity relates storage to state and control costs; it does not prove that all optimal controllers stabilize every plant under every criterion. The underlying assumptions and admissible comparisons are essential. Willems's \href{https://doi.org/10.1007/BF00276493}{Dissipative Dynamical Systems, Part I} provides the general storage-and-supply framework. Its role here is conceptual: the cost, the storage and any mechanical energy must each be identified rather than merged under one label.

\subsection{A Normalized Full-State H-Infinity Problem}

To make the disturbance claim precise, consider
\begin{equation}
\dot x=Ax+Bu+Ew,\qquad
z=\begin{bmatrix}Q^{1/2}x\\R^{1/2}u\end{bmatrix},
\quad Q\succeq0,\ R\succ0.
\end{equation}
This normalized full-state setup includes control effort in the regulated output and has no additional feedthrough or cross terms. More general output-feedback plants require different formulas and additional solvability conditions. For a chosen $\gamma>0$, seek an appropriate stabilizing symmetric solution of the game Riccati equation
\begin{equation}
A^TS+SA+Q-SBR^{-1}B^TS+
\gamma^{-2}SEE^TS=0.
\end{equation}
Choosing any positive matrix root without checking the stabilizing branch is insufficient. In the usual game formulation the associated saddle dynamics involve $A-(BR^{-1}B^T-\gamma^{-2}EE^T)S$; the existence and admissibility conditions must hold. With $u=-R^{-1}B^TSx$, completing the disturbance square yields
\begin{equation}
\dot V+z^Tz-\gamma^2w^Tw
=-\gamma^2\|w-\gamma^{-2}E^TSx\|_2^2\le0.
\end{equation}
Consequently
\begin{equation}
\int_0^T z^Tz\,dt\le
\gamma^2\int_0^T w^Tw\,dt+V(x_0)-V(x(T)).
\end{equation}
For zero initial state, nonnegative storage and a well-posed internally stable closed loop, this gives the induced $L^2$-gain bound. With nonzero initial state there is an additional initial-storage contribution. This is not a pointwise requirement that $\dot V<0$ for every disturbance. A disturbance can raise storage while the integrated gain bound holds.

For example take $\dot x=-x+u+w$, $z=(x,u)^T$ and $\gamma=2$. The positive stabilizing solution satisfies $-2S+1-3S^2/4=0$, giving $S=(-2+\sqrt7)/1.5$. At $x=1$, $w=3$ and $u=-S$, the storage derivative is approximately $1.35134>0$, while $\dot V+z^Tz-4w^2\approx-33.46333$. Both facts are consistent. A smallest feasible attenuation level is generally an infimum; an exactly optimal controller at that boundary need not exist. A worst-case disturbance criterion also does not make the same controller simultaneously optimal for an unrelated expected-energy criterion.

For a golf model, identify the disturbance channel: uncertain joint torque, ground motion, shaft loading and measurement noise enter different equations. The scale of $w$ defines the meaning of $\gamma$. A numerical attenuation level with no disturbance units or output normalization cannot support a physical robustness comparison.

\section{Control and Estimation Duality}

\subsection{The Kalman--Bucy Model}

Use an explicit stochastic model,
\begin{equation}
dx=(Ax+Bu)\,dt+G_w\,dW_t,\qquad
dy=Cx\,dt+R_v^{1/2}\,dV_t,
\end{equation}
where the Wiener processes are independent, $R_v\succ0$, and $Q_w=G_wG_w^T$. With a Gaussian initial state independent of the noises, the linear Gaussian filter has
\begin{equation}
\begin{aligned}
d\hat x&=(A\hat x+Bu)\,dt+
PC^TR_v^{-1}(dy-C\hat x\,dt),\\
\dot P&=AP+PA^T+Q_w-PC^TR_v^{-1}CP.
\end{aligned}
\end{equation}
Here $P$ is an estimation-error covariance and the noise matrices are continuous-time intensities. Writing a white-noise sample as an ordinary differentiable signal obscures those units and the stochastic interpretation. A discrete sensor covariance per sample is not automatically the same matrix as a continuous-time intensity.

The algebraic filter equation is the control Riccati equation with the dual pair $(A^T,C^T)$ and corresponding weights. The time-dependent comparison additionally requires reversing time and the coefficient histories, because the control equation runs backward from a terminal condition while covariance runs forward from an initial condition. A transpose alone does not remove that boundary-value distinction. Stability and bounded covariance require appropriate detectability, stabilizability or uniform time-dependent hypotheses, not merely the existence of a written Riccati differential equation.

\subsection{Covariance and the Differential Metric}

Covariance describes uncertainty; large covariance means a direction is uncertain. If $P\succ0$, the precision $M=P^{-1}$ provides the inverse weighting of an estimation-error norm. For the deterministic homogeneous error dynamics $A_o=A-PC^TR_v^{-1}C$, differentiate the inverse to obtain
\begin{equation}
\dot M=-M\dot P M,
\qquad
\dot M+A_o^TM+MA_o=-MQ_wM-C^TR_v^{-1}C.
\end{equation}
This identity gives a precise connection between filtering and a differential stability certificate. Uniform positivity, boundedness and an appropriate strictness or observability argument are still needed for a uniform convergence rate. It does not say that a stochastic estimation error converges to zero in the continuing presence of process and measurement noise. Instead one studies its distribution, covariance, moments or the forgetting of different initial estimates under a shared observation record.

\subsection{Nonlinear Observers Need Their Full Jacobian}

For a smooth deterministic observer
\begin{equation}
\dot{\hat x}=F(\hat x,u)+L(\hat x,t)\bigl(y-h(\hat x)\bigr),
\end{equation}
the Jacobian with the measured $y$ and commanded $u$ held fixed acts on a variation $\xi$ as
\begin{equation}
J_o\xi=D_xF\,\xi-LD_xh\,\xi+
D_xL[\xi]\bigl(y-h(\hat x)\bigr).
\end{equation}
If $F=f+Gu$, then $D_xF=Df+DG[\,\cdot\,]u$. Both the input-dependent derivative and the innovation-weighted gain derivative matter. In the scalar example $F(x,u)=x^2+(1+x)u$, $h(x)=x$ and $L(x)=x^2$, at $(\hat x,u,y)=(0.3,0.7,1.2)$ the complete Jacobian is $1.75$. Omitting either derivative gives a different variational model.

A contraction calculation uses this complete $J_o$ and the metric derivative along the observer field. Defining the symmetric part as $(J+J^T)/2$ requires consistent factors of two in every rate statement. A transposed nonlinear control law is not automatically an observer, and a local extended Kalman filter covariance is not a global contraction proof. Extended Kalman filters can have conditional local convergence guarantees, but those require assumptions about observability, nonlinear remainder bounds, covariance behavior and initialization. Differential dynamic programming also contains second-order dynamics terms beyond a simple linearized Riccati step; it is not exactly the nonlinear transpose of an extended Kalman filter.

\subsection{Observability Requires a Sensor Model}

For a linearized time-varying measurement $\delta y=C(t)\delta x$ with declared positive noise weighting $R_y(t)$, one information-like finite-horizon Gramian is
\begin{equation}
W_o=\int_{t_0}^T\Phi(t,t_0)^TC(t)^T
R_y(t)^{-1}C(t)\Phi(t,t_0)\,dt.
\end{equation}
Its interpretation requires a consistent continuous measurement model and state scales; a discrete measurement sequence uses a corresponding sum with its sample covariance. A weak direction means the modeled observations poorly distinguish that initial perturbation under the specified experiment. It does not establish that a body segment is intrinsically unobservable in every task.

Multiple muscle attachments are not, by themselves, multiple independent measurements. Claims that a shoulder is well observed while a wrist is poorly observed require a sensor map, retained states, noise assumptions and an experiment. Motion capture, inertial sensors, force plates and shaft measurements constrain different combinations of configuration, velocity and loading. Correlated sensor errors and contact modes can alter the information substantially. State estimation is therefore part of the mechanical argument, not a generic assurance that all required joint states are available.

\section{Separation, Coupling and the Swing}

\subsection{What Linear Separation Proves}

For the deterministic linear plant, use $u=-K\hat x$ and $\dot{\hat x}=A\hat x+Bu+L(y-C\hat x)$ with $y=Cx$. Writing $e=x-\hat x$ gives
\begin{equation}
\frac{d}{dt}\begin{bmatrix}x\\e\end{bmatrix}
=\begin{bmatrix}A-BK&BK\\0&A-LC\end{bmatrix}
\begin{bmatrix}x\\e\end{bmatrix}.
\end{equation}
The nominal eigenvalues are the union of the controller and observer eigenvalues. If both blocks are Hurwitz, the augmented linear system is asymptotically stable. With the linear Gaussian model and appropriate quadratic expected cost, the separation principle also supports the LQG optimality construction. This is a stronger set of hypotheses than simply combining any stable controller with any stable estimator.

Even the triangular case does not inherit the exact smaller exponential rate with a uniform prefactor in every norm. The matrix
\begin{equation}
\begin{bmatrix}-1&1\\0&-1\end{bmatrix}
\quad\hbox{has transition}\quad
e^{-t}\begin{bmatrix}1&t\\0&1\end{bmatrix}.
\end{equation}
The polynomial factor prevents a bound $Ce^{-t}$ for all $t$ with a finite constant $C$, although any strictly slower exponential rate admits a bound. Equal block decay rates and coupling can therefore matter even without nonlinearities. In a nonlinear system, bounded coupling, compatible metrics and an appropriate cascade or small-gain argument are needed; independent local contraction claims do not automatically prove stability of their interconnection. General two-way coupling can destabilize otherwise stable isolated components.

\subsection{A Technically Defensible Golf Argument}

The central opportunity is to connect mechanics, control authority and information along a feasible swing. Begin with retained states and input ports, including the shaft and actuator dynamics relevant to the question. Resolve ground and hand constraints, and declare the mode before differentiating. Compute the variational dynamics along the proposed trajectory, then map those perturbations to specified impact outputs. This determines what a local control calculation is trying to correct.

Next identify what is measured and when. An input can be mechanically effective but unavailable within the remaining activation or computation delay. A mechanically important state direction can be hard to estimate. An apparently redundant joint motion can preserve the first-order impact output while changing effort, contact forces or future sensitivity. The Riccati value expresses a selected tradeoff among such consequences; it does not establish a unique biological objective from one observed swing.

Test the proposed policy against finite perturbations in the nonlinear constrained model, with actuator limits and estimation errors included. Report the region, horizon, output changes and uncertainty under which the conclusions hold. Distinguish reducing impact dispersion from driving every joint state toward one reference: several state trajectories can be acceptable if they produce the same task outcome within physical constraints. A universal instruction to make every segment maximally stiff or every state error vanish does not follow from LQR.

Finally compare alternative hypotheses using observable predictions. Changing a weighting matrix can alter a predicted correction strategy without changing the mechanics; changing the plant can alter reachable outcomes even with fixed weights. If multiple objectives explain the same kinematics, additional measurements or perturbation experiments are needed. Technical rigor strengthens the argument by revealing which conclusion comes from mechanics, which from the chosen objective, and which still needs evidence about the golfer.

\section{Reproducible Counterexamples}

The following NumPy/SciPy program checks the scalar rates, the uncontrollable but well-conditioned value, the curved optimal path, the observer cascade, the disturbance balance, the full observer derivative and the saturation failure. Its dimensionless examples isolate mathematical claims; they are not fitted physiological models. The worked discrete table above is separately generated from its declared matrices and checked against an independent algebraic Riccati solver.

\begin{lstlisting}[language=Python]
# Duality chapter checks
import numpy as np
from scipy.linalg import expm, solve_continuous_are, solve_discrete_are


def duality_cases() -> dict[str, object]:
    """Dimensionless counterexamples; these are not fitted golf models."""
    one = np.ones((1, 1))
    scalar_s = float(solve_continuous_are(one, one, one, one).item())
    scalar_closed = 1.0 - scalar_s
    scalar_value_rate = (1.0 + scalar_s**2) / scalar_s
    discrete_a = np.array([[1.2]])
    discrete_s = solve_discrete_are(discrete_a, one, one, one)
    discrete_gain = float((discrete_s @ discrete_a / (one + discrete_s)).item())
    discrete_closed = 1.2 - discrete_gain
    discrete_value_factor = 1.0 - (1.0 + discrete_gain**2) / discrete_s.item()
    stable_s = solve_continuous_are(-np.eye(2), np.zeros((2, 1)), np.eye(2), one)

    # Fully actuated, but its optimal state path is a parabola rather than a line.
    path_a = np.diag([-1.0, -2.0])
    path_s = solve_continuous_are(path_a, np.eye(2), np.diag([3.0, 12.0]), np.eye(2))
    path_closed = path_a - path_s
    velocity = path_closed @ np.ones(2)
    acceleration = path_closed @ velocity

    # Two decay rates of one, coupled in a triangular observer/controller cascade.
    cascade = expm(np.array([[-1.0, 1.0], [0.0, -1.0]]) * 4.0)

    # xdot=-x+u+w, z=(x,u), gamma=2; solve the scalar game Riccati equation.
    gamma = 2.0
    game_s = (-2.0 + np.sqrt(7.0)) / 1.5
    state, disturbance = 1.0, 3.0
    control = -game_s * state
    storage_rate = 2 * game_s * state * (-state + control + disturbance)
    balance = storage_rate + state**2 + control**2 - gamma**2 * disturbance**2

    # F=x**2+(1+x)u, h=x, L=x**2, with held u and measured y.
    estimate, command, measurement = 0.3, 0.7, 1.2
    observer_jacobian = (
        2 * estimate + command + 2 * estimate * (measurement - estimate) - estimate**2
    )
    return {
        "scalar_norm_rate": scalar_value_rate / 2,
        "scalar_value_rate": scalar_value_rate,
        "scalar_transition": np.exp(scalar_closed * 0.7),
        "discrete_closed_loop": discrete_closed,
        "discrete_value_factor": discrete_value_factor,
        "discrete_norm_factor": np.sqrt(discrete_value_factor),
        "uncontrolled_riccati": stable_s,
        "path_velocity": velocity,
        "path_acceleration": acceleration,
        "cascade_transition": cascade,
        "disturbed_storage_rate": storage_rate,
        "hinf_balance": balance,
        "observer_jacobian": observer_jacobian,
        "saturated_state_rate": 2.0 - min(2 * scalar_s, 1.0),
        "unsaturated_state_rate": 2.0 - 2 * scalar_s,
    }


if __name__ == "__main__":
    for name, result in duality_cases().items():
        print(name, result)
\end{lstlisting}

\section{Summary and Further Reading}

The Riccati completion identities make stability and optimality meet in a specific quadratic problem. Their exact content includes the input-effort term, the correct continuous or discrete equation, the boundary condition, and the distinction between value and norm rates. Positive definiteness, uniformity and admissibility are assumptions to establish, not consequences of notation.

Geometry provides norms and differential certificates when the relevant matrices are positive and the coordinate construction is valid. It does not turn every optimal swing into a geodesic. Robustness requires the actual uncertainty channel and implementation; estimation requires a measurement model; and finite-horizon swing performance requires an impact objective and remaining control authority. These connections are useful precisely because their limits can be tested.

The continuous, finite-horizon and discrete LQR derivations can be compared with Russ Tedrake's \href{https://underactuated.mit.edu/lqr.html}{Linear Quadratic Regulators}. The full-state and estimator constructions and return-difference analysis appear in Alexandre Megretski's \href{https://ocw.mit.edu/courses/6-245-multivariable-control-systems-spring-2004/b49e52078cd6ec5648ce2d66dff869f2_lec4_6245_2004.pdf}{Using H2 Optimization}. Use the explicit negative-feedback algebra in this chapter when comparing sign conventions. Megretski's \href{https://ocw.mit.edu/courses/6-245-multivariable-control-systems-spring-2004/a9100bcb2c4fe5f05ef2d86007d1fd65_lec7_6245_2004.pdf}{Algorithms for H-Infinity Optimization} states the additional solvability conditions for more general plants; Tedrake's \href{https://underactuated.mit.edu/robust.html}{Robust and Stochastic Control} develops the normalized disturbance-attenuation setting. Bishop and Del Moral's \href{https://arxiv.org/abs/1610.04686}{On the Stability of Kalman--Bucy Diffusion Processes} treats the stochastic filter and conditional stability questions. These sources support particular mathematical frameworks, not an empirical identification of the golfer's objective.

\section{Exercises}

\begin{enumerate}
\item Derive the continuous completion identity without assuming $u=-Kx$ until the final step. Keep the finite terminal value and specify when the infinite-horizon limit is legitimate. Construct a stable unpenalized state for which the Riccati value is singular.

\item Expand the discrete Bellman identity and recover both $Q$ and $K^TRK$ in the value difference. For $A=1.2$ and $B=Q=R=1$, calculate the state and value multipliers and convert the state multiplier to a decay rate for a declared sample period.

\item Recompute every row of the generated example. Compare the actual value ratio, the conservative bound and the spectral radius. Explain the initial Euclidean norm growth without claiming instability. Then derive a consistent exact zero-order-hold pendulum model and quantify how its LQR solution differs.

\item Solve a finite-horizon problem with zero terminal cost. Show why a positive-definite metric bound fails at the endpoint, and identify additional terminal ingredients needed for a proposed receding-horizon stability argument.

\item Rescale one coordinate in the uncontrollable stable example by a factor of one thousand. Transform all matrices consistently, compare the two Riccati condition numbers, and show that the physical trajectories and control authority are unchanged.

\item Verify the curved optimal path in the constant-metric example and calculate its velocity and acceleration determinant. Explain why neither an optimal-control cost nor a quadratic value is automatically the length of that state curve.

\item Derive the common scalar gain certificate and the SISO return-difference identity. For the scalar unstable plant, compare the guaranteed gain interval with the exact stability interval. Explain why the Nyquist disk is not a pole-location disk and why neither result covers actuator saturation.

\item Complete the disturbance square for the normalized H-infinity problem. Reproduce the positive storage derivative example and the negative supply balance. State the initial-state, disturbance-unit and internal-stability assumptions needed to interpret an induced gain.

\item Derive the precision-metric identity from the Kalman--Bucy covariance equation. Distinguish decay of a homogeneous error from nonzero steady-state covariance under persistent noise. Explain the time reversal needed for finite-horizon control/filter Riccati duality.

\item Differentiate the state-dependent observer with held measurements and input, including the gain derivative. Verify the scalar Jacobian by finite differences. Identify the additional assumptions required to turn a pointwise calculation into a regional contraction claim.

\item Propose a swing experiment with declared retained states, actuator port, sensors, noise scales, impact outputs and contact assumptions. Identify one perturbation that is costly to correct and one that is difficult to observe. Explain which measurements could distinguish a change in mechanics from a change in the assumed objective, and which conclusions remain model-dependent.
\end{enumerate}
